\section*{Verificación y Mantenimiento}

\begin{itemize}
  \item Verificar que la fuente regreso al contenedor (con el Geiger).
  \item Retirar señales y acotamientos.
  \item No arrastrar mangueras y desensamblar el equipo al moverse a otra área.
\end{itemize}

\section*{Acordonamiento y señalización de áreas.}
\textbf{Ya cubierto.}

\section*{Uso y cuidado de los dosímetros de lectura directa y de película}
\begin{itemize}
  \item Se usarán siempre que se esté trabajando con radiación.
  \item Se usarán al frente y entre el cuello y la cintura.
  \item Se vigilará frecuentemente el dosímetro de lectura directa y se recargará cuando llegue a 100 mR o antes.
  \item No se debe humedecer el dosímetro de película.
  \item No se debe almacenar el dosímetro con la fuente.
\end{itemize}

\section*{Operación y mantenimiento del equipo de gammagrafía}
Efectuarlo cada dos meses o antes si es necesario.

\section*{Operación y cuidado de Detectores de Radiación y Alarma}
\begin{itemize}
  \item Evitar la humedad.
  \item Evitar golpes y maltratos.
  \item Cuidarlos de polvo, aceite, grasa, etc.
  \item Almacenarlos cuando no se estén usando.
\end{itemize}

\section*{Calibración de Dosímetros de lectura directa}
Solo administrativo.

\section*{Calibración del Detector Geiger Victoreen 492}
Solo administrativo.

\chapter{Criterios a seguir en emergencias}
\textbf{Garantizar la seguridad física de usted y acompañantes.}
\begin{itemize}
  \item Garantizar la integridad del contenedor y la fuente.
  \item Proteger al público.
\end{itemize}

Siempre se debe llenar la forma SR-12, ``Reporte de Emergencia Radiológica''.

\textit{Se debe consultar el Manual para ver las instrucciones específicas de cada accidente.}
